\documentclass[11pt, a4paper]{article}

% --- Packages ---
\usepackage[utf8]{inputenc}
\usepackage[T1]{fontenc}
\usepackage{geometry}
\geometry{top=2.5cm, bottom=2.5cm, left=2.5cm, right=2.5cm}
\usepackage{mathpazo} % Palatino font
\usepackage[dvipsnames]{xcolor}
\usepackage{hyperref}
\usepackage{listings}
\usepackage{tcolorbox}
\usepackage{graphicx}
\usepackage{enumitem}
\usepackage{fancyhdr}
\usepackage{booktabs}

% --- Styling ---
\definecolor{codegray}{rgb}{0.95,0.95,0.95}
\definecolor{codeblue}{rgb}{0.1,0.1,0.5}
\definecolor{codegreen}{rgb}{0,0.5,0}

\lstset{
    backgroundcolor=\color{codegray},
    commentstyle=\color{codegreen},
    keywordstyle=\color{magenta},
    numberstyle=\tiny\color{gray},
    stringstyle=\color{codeblue},
    basicstyle=\ttfamily\small,
    breaklines=true,
    numbers=left,
    numbersep=5pt,
    frame=single
}

\pagestyle{fancy}
\fancyhf{}
\rhead{\textbf{ClusterLLM-Local}}
\lhead{Day 2: System Reproducibility}
\rfoot{\thepage}

% --- Document Info ---
\title{\textbf{\Huge Day 2 Technical Report}\\
\Large System-Level Reproducibility with Docker \& GPU Containers}
\author{\textbf{Hamady Gackou} \\ Master Data Science \& AI}
\date{\today}

\begin{document}

\maketitle
\thispagestyle{empty}

\begin{abstract}
This document reports the engineering work completed on Day 2 of the \textit{ClusterLLM-Local} project.
While Day 1 established reproducibility at the \textbf{code} and \textbf{data} levels using Poetry and DVC,
Day 2 addresses a deeper and often overlooked layer: \textbf{system-level reproducibility}.
We containerize the entire execution environment using Docker and Docker Compose, including the operating system,
CUDA runtime, Python interpreter, and dependency installation process.
The result is a deterministic, GPU-enabled environment that can be rebuilt and executed on any compatible machine
without manual configuration.
\end{abstract}

\tableofcontents
\vspace{1cm}
\hrule
\vspace{1cm}

% =========================================================
\section{From Day 1 to Day 2: Why Containerization Is Necessary}

\subsection{What Day 1 Did Not Solve}
Day 1 successfully locked:
\begin{itemize}
    \item Python dependencies via \texttt{poetry.lock}
    \item Data versions via \texttt{DVC}
\end{itemize}

However, the project still depended on:
\begin{itemize}
    \item the host operating system
    \item system libraries (libc, SSL, zlib, etc.)
    \item CUDA runtime versions
    \item Python interpreter build
\end{itemize}

These elements are \textbf{outside the control of Poetry and DVC}.

\subsection{The Remaining Risk}
Even with perfect code and data locking:
\begin{itemize}
    \item a different Linux distribution
    \item a different CUDA version
    \item a missing system library
\end{itemize}
can silently break the project or change results.

\begin{tcolorbox}[colback=red!5!white,colframe=red!75!black,title=Key Observation]
Reproducibility is not achieved until the \textbf{entire system stack} is controlled.
\end{tcolorbox}

% =========================================================
\section{System-Level Reproducibility: Definition}

\subsection{Formal Definition}
System-level reproducibility is the ability for an independent researcher to reconstruct:
\begin{itemize}
    \item the same operating system
    \item the same CUDA runtime
    \item the same Python interpreter
    \item the same dependency installation procedure
\end{itemize}
using only declarative configuration files.

\subsection{Why It Matters for This Project}
\begin{itemize}
    \item \textbf{GPU Dependence:} Results depend on CUDA and cuDNN behavior.
    \item \textbf{Academic Evaluation:} A jury must be able to rerun the system.
    \item \textbf{Longevity:} The project must remain executable months later.
\end{itemize}

% =========================================================
\section{Why Docker?}

\subsection{Rejected Alternatives}
\begin{itemize}
    \item \textbf{Manual Setup:} Non-reproducible, undocumented, error-prone.
    \item \textbf{Virtual Machines:} Heavy, non-versioned, hard to audit.
\end{itemize}

\subsection{Docker as Infrastructure-as-Code}
Docker allows the entire system to be described as code:
\begin{itemize}
    \item versioned
    \item reviewable
    \item reproducible
\end{itemize}

\begin{tcolorbox}[colback=green!5!white,colframe=green!75!black,title=Analogy]
\textbf{Poetry} locks the ingredients.\\
\textbf{Docker} locks the kitchen, the oven, and the gas pressure.
\end{tcolorbox}

% =========================================================
\section{Dockerfile Design}

\subsection{Base Image Selection}
\begin{lstlisting}[language=docker]
FROM nvidia/cuda:12.4.1-cudnn-runtime-ubuntu22.04
\end{lstlisting}

\textbf{Rationale:}
\begin{itemize}
    \item Official NVIDIA-maintained image
    \item Explicit CUDA and cuDNN versions
    \item Compatibility with NVIDIA L4 GPU
\end{itemize}

\subsection{Explicit System Dependencies}
System libraries are installed explicitly to avoid hidden assumptions:
\begin{lstlisting}[language=docker]
RUN apt-get update && apt-get install -y \
    build-essential wget curl git ca-certificates \
    libssl-dev zlib1g-dev libbz2-dev libreadline-dev \
    libsqlite3-dev libffi-dev liblzma-dev
\end{lstlisting}

\subsection{Python 3.13 Built from Source}
Ubuntu 22.04 does not ship Python 3.13. Therefore, Python is compiled manually:
\begin{lstlisting}[language=docker]
wget Python-3.13.0.tgz
./configure --enable-optimizations
make altinstall
\end{lstlisting}

This ensures:
\begin{itemize}
    \item full control of interpreter version
    \item independence from OS package repositories
\end{itemize}

\subsection{Ensuring pip Availability}
\begin{lstlisting}[language=docker]
python3.13 -m ensurepip --upgrade
\end{lstlisting}

This step is critical: Python compiled from source does not always include \texttt{pip}.

% =========================================================
\section{Poetry Inside the Container}

\subsection{Why Poetry Must Live Inside Docker}
Installing Poetry on the host is insufficient.
The container must be fully autonomous.

\begin{lstlisting}[language=docker]
python3.13 -m pip install poetry
\end{lstlisting}

\subsection{Deterministic Dependency Installation}
\begin{lstlisting}[language=docker]
COPY pyproject.toml poetry.lock ./
RUN poetry install --no-root
\end{lstlisting}

The lock file becomes the \textbf{single source of truth} for the Python environment.

% =========================================================
\section{Docker Compose: Execution and Orchestration}

\subsection{Why Docker Compose}
Docker builds images.
Docker Compose defines how they run.

Compose controls:
\begin{itemize}
    \item GPU access
    \item volumes
    \item working directories
    \item interactive execution
\end{itemize}

\subsection{GPU Exposure}
\begin{lstlisting}[language=yaml]
gpus: all
\end{lstlisting}

This connects the container runtime to the NVIDIA driver through the container toolkit.

\subsection{Volume Mounts}
\begin{lstlisting}[language=yaml]
volumes:
  - ../:/workspace
  - ~/.cache/huggingface:/root/.cache/huggingface
  - ~/.cache/dvc:/root/.cache/dvc
\end{lstlisting}

These mounts:
\begin{itemize}
    \item avoid repeated downloads
    \item separate compute from storage
    \item improve iteration speed
\end{itemize}

% =========================================================
\section{Implementation Log (Day 2)}

\subsection{Build}
\begin{lstlisting}[language=bash]
docker compose -f docker/docker-compose.yml build --no-cache
\end{lstlisting}

Forcing a clean build ensures no hidden dependencies exist.

\subsection{Run}
\begin{lstlisting}[language=bash]
docker compose -f docker/docker-compose.yml run clusterllm
\end{lstlisting}

Successful execution validates:
\begin{itemize}
    \item image integrity
    \item CUDA availability
    \item correct system configuration
\end{itemize}

% =========================================================
\section{Critical Evaluation}

\subsection{What Day 2 Solves}
\begin{itemize}
    \item OS-level reproducibility
    \item CUDA reproducibility
    \item Python interpreter determinism
    \item elimination of “works on my machine”
\end{itemize}

\subsection{Trade-offs}
\begin{itemize}
    \item long build times (Python compilation)
    \item larger image size
    \item increased infrastructure complexity
\end{itemize}

These costs are intentional and justified by scientific rigor.

% =========================================================
\section{Conclusion and Transition}

Day 2 completes the transition from:
\begin{quote}
Reproducible code and data
\end{quote}
to:
\begin{quote}
Reproducible system, code, and data
\end{quote}

\textbf{Next Step (Day 3):} Experimental pipelines, embeddings, clustering strategies, and evaluation metrics.

\end{document}
